\documentclass{article}
\usepackage{shortex}
\usepackage[margin=1in]{geometry}
\usepackage{listings}
\usepackage{xcolor}
% \usepackage{jupynotex}
\usepackage{csvsimple}
\usepackage{placeins}
% \usepackage{hyperref}


\graphicspath{{./graphics}}

\lstdefinelanguage{Julia}{
    keywords={function, end, if, else, elseif, for, while, return, 
               local, global, const, struct, mutable, abstract, type,
               using, import, export, module, begin, let, do, try, 
               catch, finally, true, false, nothing},
    sensitive=true,
    comment=[l]{\#},
    morecomment=[s]{\#=}{=\#},
    morestring=[b]",
    morestring=[b]',
}

\lstset{
    language=Julia,
    basicstyle=\ttfamily\small,
    keywordstyle=\color{blue}\bfseries,
    commentstyle=\color{gray}\itshape,
    stringstyle=\color{purple},
    numberstyle=\tiny\color{gray},
    numbers=left,
    stepnumber=1,
    frame=single,
    breaklines=true,
    captionpos=b,
    tabsize=4,
    showstringspaces=false,
}

\lstset{
    literate={ρ}{$\rho$}{1}
             {α}{$\alpha$}{1}
             {β}{$\beta$}{1}
             {σ}{$\sigma$}{1}
             {σ²}{$\sigma^2$}{2}
             {μ}{$\mu$}{1}
             {ϵ}{$\epsilon$}{1}
             {η}{$\eta$}{1}
             {η²}{{$\eta^2$}}{2}
             {r̂}{$\hat r$}{1}
             {Δ}{$\Delta$}{1}
}

\begin{document}
\title{STAT 547E - Homework 1}
\author{Zachary Lau}
\date{\today}
\maketitle

\section{Problem 1}
\begin{enumerate}
    \item % a)
    We begin by proving the hint. Firstly we show that under ELE
    $\E g^n_t = \frac {Z_{\beta_t}}{Z_{\beta_{t-1}}}$
    By the ELE assumption we can replace expectations in
    $g(X_{t-1})$ with expectations over $X_{t-1}$ \emph{even if $X_{t-1}$ is not
    actually distributed according to $\pi_{\beta_{t-1}}$}.
    \begin{align}
        \E g^n_t &= \E_{\pi_{\beta_{t-1}}}
            \frac{\gamma_{\beta_{t}}(X_{t-1})}{\gamma_{\beta_{t-1}}(X_{t-1})} \\
        &= \int \pi_{\beta_{t-1}}(\d x_{t-1})
            \frac{\gamma_{\beta_{t}}(x_{t-1})}{\gamma_{\beta_{t-1}}(x_{t-1})} \\
        &= \int \pi_{\beta_{t-1}}(\d x_{t-1})
            \frac{\pi_{\beta_{t}}(x_{t-1}) Z_t}{\pi_{\beta_{t-1}}(x_{t-1}) Z_{t-1}} \\ 
        &= \frac{Z_t}{Z_{t-1}} 
    \end{align}
    In the notation of the question we get $z_t := \frac{Z_t}{Z_{t-1}}$, and
    therefore the normalized weights have expectation $1$.
    We now consider the \emph{variance} of the normalized weights $\frac{g_n^t}{z_t}$.
    Playing the same trick with expectations under ELE we get
    \begin{align}
        \Var g^n_t &= \E \left[ \left(\frac{g^n_t}{z_t} - 1\right)^2 \right] \\
        &= \E_{\pi_{\beta_{t-1}}}
        \left[ \left(\frac{\gamma_{\beta_{t}}(X_{t-1})/Z_t}
        {\gamma_{\beta_{t-1}}(X_{t-1})/Z_{t-1}} - 1\right)^2 \right] \\
        &=  \E_{\pi_{\beta_{t-1}}}
        \left[ \left(\frac{\pi_{\beta_{t}}(X_{t-1})}
        {\pi_{\beta_{t-1}}(X_{t-1})} - 1\right)^2 \right] \\
        &= \chi^2(\pi_{\beta_t} \| \pi_{\beta_{t-1}})
    \end{align}

    Now that the hint is proven we may tackle the variance of the estimators.
    We begin with AIS. Each particle is independent, so we just have to consider
    a single particle going from start to finish. We begin by showing that the
    AIS estimator is unbiased. We use ELE to get independence and split
    up the expectation of the product as a product of expectations. This leaves
    a telescoping product which can be reduced in terms of its first and
    last terms.

    \newcommand{\AIS}{{\text{AIS}}}
    \begin{align}
        \E \hat Z^n_{\AIS} &= 
            \E \prod_{t=1}^T g^n_T \\
        &= \prod_{t=1}^T \E g^n_T \\
        &= \prod_{t=1^T} \frac{Z_T}{Z_{T-1}} \\
        &= \frac{Z_T}{Z_0}
    \end{align}
    Assuming the refernece is normalized and that $Z_T := Z$ as in the
    question, we see that for each particle the estimator $\hat Z^n_{\AIS}$
    is unbiased. As an arithmetic mean of unbiased estimators, the overall 
    estimator $\hat Z_\AIS$ is also unbiased.
    
    We now tackle the variance, again
    independence lets us break up products of exepctations as expectations
    of products. Further breaking up $Z = \prod_{t=1}^T z_t$ we get
    \begin{align}
        \E \left(\frac{{\hat Z}^{n}_\AIS}{Z}\right)^2 &=
            \E \left(\prod_{t=1}^T \frac{g^n_t}{z_t} \right)^2 \\
            &= \prod_{t=1}^T \left(\E {\frac{g_t^{n}}{z_t}}\right)^2 \\
            &= \prod_{t=1}^T \Var \frac{g^n_t}{z_t} + \left(\E \frac{g^n_t}{z_t}\right)^2 \\
            &= \prod_{t=1}^T \chi^2(\pi_{\beta_t}\| \pi_{\beta_{t-1}}) + 1 
    \end{align}
    Recognizing that each particle's path is independent, and letting
    $D_t = \log\left(1+\chi^2(\pi_{\beta_t}\|\pi_{\beta_{t-1}})\right)$ we
    see that the variance of the overall estimator is given by 
    \begin{align}
        \Var \left(\frac{{\hat Z}^{n}_\AIS}{Z}\right) &=
           \frac{\prod_{t=1}^T \chi^2(\pi_{\beta_t}\| \pi_{\beta_{t-1}}) + 1 }{N}  \\
           &= \frac{\exp\left(\sum_{t=1}^T D_t\right) - 1}{N}
    \end{align}
    
    \newcommand{\SIR}{{\text{SIR}}}
    Now we derive the variance for the noramlized SIR estimator
    $\frac{\hat Z_\SIR}{Z}$. We begin by showing that it is unbiased, that
    is $\E \hat Z_\SIR = Z$

    We first introduce the notation $\hat R_t := \frac{1}{N}\sum_{n=1}^N g^n_t$
    Note that $\hat Z_\SIR = \prod_{t=1}^T R_t$, and furthermore under ELE the
    $R_t$ are independent for different $t$. We thus turn our attention to
    computing the expectation and variance of each individual $\hat R_t$. Again
    under ELE the individual $g^n_t$ are independent across $n$, and so the
    previous results apply
    \begin{align}
        \E \hat R_t &= \E g^n_t  \\
            &= z_t = \frac{Z_t}{Z_{t-1}}
    \end{align}
    And
    \begin{align}
        \Var \left(\frac{\hat R_t}{z_t}\right) &=
            \frac 1 N \Var \left( \frac{g^n_t}{z_t}\right) \\
            &= \frac {\chi^2(\pi_{\beta_t} \| \pi_{\beta_{t-1}})}{N}  \\
            &= \frac{\exp(D_t) - 1}{N}
    \end{align}
    Given this, we can approach the overall estimator
    \begin{align}
        \E \hat Z_\SIR &= \E \prod_{t=1}^T R_t \\
        &= \prod_{t=1}^T \E R_t \\
        &= \prod_{t=1}^T z_t \\
        &= \frac{Z_T}{Z_0} \\
        &= Z
    \end{align},
    again assuming a normalized reference and ELE.
    \begin{align}
        \E\left( \frac{\hat Z_\SIR}{Z} \right)^2 &= 
            \prod_{t=1}^T \Var \left( \frac{\hat R_t}{z_t} \right) + 1 \\
            &= \prod_{t=1}^T \left(
               \frac{\exp(D_t) - 1}{N} + 1
            \right)
    \end{align}.
    Which gives the desired result
    \begin{align}
        \Var\left(\frac{\hat Z_\SIR}{Z} \right) &= 
            \prod_{t=1}^T \left(
               \frac{\exp(D_t) - 1}{N} + 1
            \right) - 1 
    \end{align}.
    \item %b 
    We use a quadratic approximation of $D_t$ based on the small-signal
    quadratic approximation of the $\chi^2$-divergence. We let $u = t/T$
    repersent our normalized index. Denote
    $f(u') = \chi^2(\pi_{\beta_{t'}} \| \pi_{\beta_t})$, 
    $\beta_{t'} = \varphi(u')$ and $\beta_t = \varphi(u)$,
    and compute a Taylor
    approximation around $u'= u$. First we evaluate
    \[
        f(u) = \chi^2(\pi_{\beta_t} \| \pi_{\beta_t}) = 0
    \]
    Taking derivatives and assuming sufficient regularity conditions are
    satifies such that expectations and derivatives can be exchanged we find the
    following. In the last line we use the fact that the expected value of the
    Fisher score is 0.
    \begin{align}
        f'(u') &= \frac{\d}{\d u} \E_{\pi_{\beta_u}} \left[ \left(
            \frac{\pi_{\beta_{u'}}}{\pi_{\beta_{u}}} - 1
            \right)^2 \right] \\
        &= \frac{\d}{\d u}\left( \E \left[ \left(
            \frac{\pi_{\beta_{u'}}}{\pi_{\beta_{u}}}
            \right)^2  - 1\right]\right) \\
        &= \frac{\d}{\d t}\left( \E \left[ \left(
            \frac{\pi_{\beta_{u'}}}{\pi_{\beta_{u}}}
            \right)^2 \right] \right) \\ 
        &=   \frac{\d}{\d t}\E \left[ \left(
            \frac{\pi_{\beta_u'}}{\pi_{\beta_{u}}}
            \right)^2 \right] \\
        &= \E\left[2\left(
            \frac{\pi_{\beta_u'}}{\pi_{\beta_{u}}}
            \right)\left(\frac 1 {\pi_{\beta_u}}\right)
            \left(\frac{\d \pi_{\beta_{u'}}}{\d \beta_{t'}}\right)
            \frac{\d \beta_{u'}}{\d t'}\right] \\
        &= 2\frac{\d \beta_{u'}}{\d t'}
            \E\left[
                \left(
            \frac{\pi_{\beta_u'}}{\pi_{\beta_{u}}}
            \right)\left(\frac 1 {\pi_{\beta_u}}\right)
            \left(\frac{\d \pi_{\beta_{u'}}}{\d \beta_{t'}}\right)
            \right] \\
        f'(u) &= 
            2\frac{\d \beta_{u'}}{\d t'}
            \E\left[
            \left(\frac 1 {\pi_{\beta_u}}\right)
            \left(\frac{\d \pi_{\beta_u}}{\d \beta_u}\right)
            \right]  \\
            &= 2\frac{\d \beta_{u'}}{\d t'}\E\left[\frac{\d}{\d\beta_u}\log \pi_{\beta_t}\right] \\
            &= 0
    \end{align}
    Taking another derivative with the product rule yields
    \begin{align}
        f''(u') &= 
            2\left(\frac{\d \beta_{u'}}{\d u'}\right)^2 \E
            \left[
                \left(
                    \frac 1 {\pi_{\beta_u}^2}
                \right)
                \left(
                    \frac {\d \pi_{\beta_{u'}}}{\d \beta_{u'}}
                \right)^2
                + 
                \left(\frac{\pi_{\beta_{u'}}}{\pi_{\beta_u}^2}\right)
                \left(\frac{\d^2 \pi_{\beta_{u'}}}{\d \beta_{u'}}\right)
            \right]
    \end{align}
    Note that similar to the Fisher score, the second summand has expectation
    zero at $\beta_{t'} = \beta_t$
    \begin{align}
        \int \pi_{\beta_u}\left(\frac{\pi_{\beta_{u}}}{\pi_{\beta_t}^2}\right)
        \frac{\d^2 \pi_{\beta_{u}}}{\d \beta_{u}} &= 
            \frac{\d^2}{\d u^2} \int \pi_{\beta_u} \\
        &= \frac{\d^2}{\d u^2} 1 \\
        &= 0
    \end{align}
    Therefore we get
    \begin{align}
        f''(u) &= 
            2\left(\frac{\d \beta_{u}}{\d u}\right)^2 \E \left[
                \left(
                    \frac 1 {\pi_{\beta_u}^2}
                \right)
                \left(
                    \frac {\d \pi_{\beta_{u}}}{\d \beta_{u}}
                \right)^2 
            \right] \\
            &=  2\left(\frac{\d \beta_{u}}{\d u}\right)^2 I(\beta_u)
    \end{align}
    Combining these derivative and recalling that
    $\frac{d}{\d u}\beta_u = \dot\varphi(u)$ gives small scale approximation
    \[
        \chi^2 (\pi_{\beta_{u'}} \| \pi_{\beta_u}) 
        = \dot\varphi(u)^2 I(\beta_u)(u'-u)^2 + o((u'- u)^3)
    \]
    For our fixed $t$ schedule, we find $u'-u = \frac 1 T$ and so we get
     \[
        \chi^2 (\pi_{\beta_{u'}} \| \pi_{\beta_u}) 
        = \dot\varphi(t)^2 I(\beta_u)T^{-2} + o(T^{-3})
    \]

    Furthermore we use the following small-scale functional approximations near
    $x = 0$ to
    get an overall approximation of the variance terms
    \begin{align}
        \log(1+x) &= x + o(x^2) \\
        \exp(x)-1 &= x + o(x^2)
    \end{align}

    Thus we find
    \[
        D_t = \dot\varphi\left(\frac t T\right)^2 I(\beta_t)T^{-2} + o(T^{-3})
    \]
    Assuming appropriate regularity conditions (e.g. uniform continuity of
    $I(\beta)$)
    \begin{align}
        \sum_{t=1}^T D_t &= 
            \frac 1 T \sum_{t=1}^T \dot\varphi(u)^2 I(\varphi(u))\frac 1 T + o(T^{-2}) \\
            &= \frac 1 T \int_{0}^{1} I(\varphi(u))\dot\varphi(u)^2 \d u + o(T^{-2})
    \end{align}
    This gives 
    \begin{align}
        \exp\left(\sum_{t=1}^T D_t \right) - 1 &=
            \frac 1 T \int_{0}^{1} I(\varphi(u))\dot\varphi(u)^2 \d u + o(T^{-2}) 
    \end{align}
    And hence building on previous results
    \begin{align}
        T\Var \left[\frac{\hat Z_\AIS}{Z}\right] &= 
            \frac 1 N \int_{0}^{1} I(\varphi(u))\dot\varphi(u)^2 \d u + o(N^{-1}T^{-1})
    \end{align}
    Which in the limiting case with fixed $N$ gives
    \[
        \lim_{T \to \infty} T\Var \left[\frac{\hat Z_\AIS}{Z}\right] = 
            \frac 1 N \int_{0}^{1} I(\varphi(u))\dot\varphi(u)^2 \d u = 
            \frac{E[\varphi]}{N}
    \]

    For SIR the product can be turned into a sum using logarithms hence
    \begin{align}
        \prod_{t=1}^T \left(1 + \frac{\exp(D_t)-1}{N}\right) - 1
        &= \exp\left(\sum_{t=1}^T 
            \log\left(1 + \frac{\exp(D_t)-1}{N} \right)
            \right) - 1
    \end{align}
    Using our above small-scale approximations we get
    \[
        \exp(D_t) - 1 =  \dot\varphi\left(\frac t T\right)^2 I(\beta_t)T^{-2} + o(T^{-3})
    \]
    And 
    \[
        \log\left( 1 + \frac{\dot\varphi\left(\frac t T\right)^2 I(\beta_t)T^{-2} + o(T^{-3})}{N}\right) = 
        \frac 1 N\dot\varphi\left(\frac t T\right)^2 I(\beta_t)T^{-2} + o(N^{-1}T^{-3})
    \]
    Then
    \[
        \sum_{t=1}^T
             \left[ \frac 1 N\dot\varphi\left(\frac t T\right)^2 I(\beta_t)T^{-2} + o(N^{-1}T^{-3})\right]
            =  \frac 1 {NT} \int_0^1 I(\varphi(u))\dot\varphi(u)^2 \d u
                + o(N^{-1}T^{-2})
    \]
    So pulling it all together with the $\exp$ identity
    \[
        \prod_{t=1}^T \left(1 + \frac{\exp(D_t)-1}{N}\right) - 1 =
         \frac 1 {NT} \int_0^1 I(\varphi(u))\dot\varphi(u)^2 \d u
                + o(N^{-1}T^{-2}) 
    \]
    Again, using this to analyze the variance of SIR we get
    \[
        T\Var \left[ \frac{\hat Z_\SIR}{Z} \right] = 
         \frac 1 {N} \int_0^1 I(\varphi(u))\dot\varphi(u)^2 \d u
                + o(N^{-1}T^{-1})
    \]
    or in the asymptotic case of fixed $N$
    \[
        \lim_{T \to \infty} T\Var \left[ \frac{\hat Z_\SIR}{Z} \right] = 
         \frac 1 {N} \int_0^1 I(\varphi(u))\dot\varphi(u)^2 \d u = 
         \frac{E[\varphi]}{N} 
    \]
    
    
    % We consider a local linear approximation of $\exp(D_t)$. That is, fixing
    % $\beta_t$ we consider a linear approximation of 
    % $f(t') = 1 + \chi^2 ( \pi_{{t'}} \| \pi_{\beta_{t}})$
    % Firstly we evaluate at $t'=0$, in this case
    % $\chi^2 (\pi_{\beta_t} \| \pi_{\beta_t}) = 0$ so 
    % $f(t) = 0$. Before taking the derivative, we simplify the expression.
    % We see 
    % \begin{align}
    %     1 + \chi^2 ( \pi_{\beta_{t'}} \| \pi_{\beta_{t}}) &=
    %     1 + \E_{\pi_\beta}
    %         \left( \frac{\d \pi_{\beta_{t'}}}{\d \pi_{\beta_t}} - 1\right)^2 \\
    %     &= \E_{\pi_{\beta_t}} \left( \frac{\d \pi_{\beta_{t'}}}{\d \pi_{\beta_t}} \right)^2 
    % \end{align}
    % Taking derivatives and assuming that sufficient regularity conditions are
    % satisfied such that we can exchange expectation and differentiation, and
    % that the Fisher score exists and is finite we see
    % \begin{align}
    %     f'(t') &= \frac{\d}{\d t'}  \E_{\pi_{\beta_t}} \left( \frac{\d \pi_{{t'}}}{\d \pi_{\beta_t}} \right)^2 \\
    %     &= \E_{\pi_{\beta_t}}
    %     \left[ \frac{\d}{\d t'}  \left(\frac{\d \pi_{\beta_{t'}}}{\d \pi_{\beta_t}} \right)^2 \right] \\
    %     &= \E_{\pi_{\beta_t}}
    %     \left[
    %         2 \left(\frac{\d \pi_{\beta_{t'}}}{\d \pi_{\beta_t}}\right)
    %     \right] 
    % \end{align}

    \item % c
    Equality is immediate. Taking $\varphi(u) = F^{-1}(\Lambda u)$ and using the
    Fundamental Theorem of Calculus we get
    \begin{align}
        F(\varphi(u)) &= \Lambda u \\
        F'(\varphi(u))\dot\varphi(u) &= \Lambda \\
        \sqrt{I(\varphi(u))}\dot\varphi(u) &= \Lambda \\
        I(\varphi(u))\dot\varphi(u)^2 &= \Lambda \\
        E[\varphi] &= \int_0^1 I(\varphi(u))\dot\varphi(u)^2 \d u \\
            &= \int_0^1 \Lambda^2 \d u  \\
            &= \Lambda^2
    \end{align}
    Inequality can be shown by Jensen's. Note that using a $u$-subsitution
    \begin{align}
        \int_0^1 \sqrt{I(\varphi(u))}\dot\varphi(u) \d u &= 
            \int_0^1 \sqrt{I(\beta)}\d \beta \\
            &= \Lambda
    \end{align}
    We apply Jensen's inequality for $\psi(x)=x^2$. This is a convex function
    so we get
    \begin{align}
        \int_0^1 I(\varphi(u))\dot\varphi(u)^2 \d u &= 
            \int_0^1 \left(\sqrt{I(\varphi(u))\dot\varphi(u)}\right)^2 \d u \\
            &\geq \left[ \int_0^1 \sqrt{I(\varphi(u))\dot\varphi(u)} \d u \right]^2 \\
            &= \Lambda^2
    \end{align}
    Taken together these gives us a few insights for the design of AIS and SIR
    practice. Firstly, from results in (a) we see that the \emph{importance of
    annealing}, since without annealing the difficulty of our problem is
    exponential in the divergence between the two distributions. Results in 
    (b) confirm that adding more annealing layers reduces the variance of our
    estimators, and moreover in the large $T$ regime it does so 
    \emph{linearly in 
    $T$}. It is no, surprise, but we confirm as well that it also does so
    linearly in $N$. Thus either of these knobs can be turned to give us lower
    variance. Furthermore in the large $T$ regime both \emph{AIS and SIR are equally
    asymptotically efficient}. Finally section (c) discusses optimal annealing
    schedules for large $T$. The optimal schedule is one which has
    \emph{unit speed} in terms of Fisher-Rao distance (recalling that the Fisher
    information gives us the metric tensor).
\end{enumerate}

\section{Problem 2}
\begin{enumerate}
    \item % a
    In this case the linear path is given by 
    % We first derive the Fisher information. 
    % \begin{align}
    %     \ell(\mu) &= \log \Norm(x; \mu, \sigma^2) \\
    %         &= -\frac 1 2 \log(2\pi\sigma^2) - \frac 1 2 \frac{(x-\mu)^2}{\sigma^2} \\
    %     \frac{\partial \ell}{\partial \mu} &=
    %         -\frac{(x-\mu)}{\sigma^2} \\
    %     \frac{\partial^2 \ell}{\partial \mu^2} &= \frac{1}{\sigma^2} \\
    %     I(\mu) = \E\frac{\partial^2 \ell}{\partial \mu^2} &= \frac{1}{\sigma^2}
    % \end{align}
    % Now 
    \item %b 
    We derive the Fisher information for $\sigma^2$ fixing $\mu$ and using that
    $\E X = \mu$
    \begin{align}
        \frac{\partial \ell}{\partial \sigma^2} &=
            -\frac 1 {2\sigma^2} + \frac 1 2 \frac{(x-\mu)^2}{\sigma^4} \\
        \frac{\partial \ell^2}{{\partial \sigma^2}^2} &=
            \frac{1}{2\sigma^4} - \frac{(x-\mu)^2}{\sigma^6} \\
        I(\sigma^2) &=
            \frac{1}{2\sigma^4}
    \end{align}

\end{enumerate}
\end{document}